% GENERATED FILE. Edit canonical lessons, then run npm run generate:books.
% canonical-source-hash: fnv1a64:8b320b1cb9869013
% canonical-lessons: RU-C286-translirovat, RU-C286-osveshchat, RU-C286-raskryvat, RU-C286-publikatsiya, RU-C286-tirazh, RU-R286-review-279-286

\chapter{Covering a story}
\label{ch:ru-covering-a-story}
\hlchaptermodality{286}

\begin{chapteropening}
\textbf{By the end of this chapter:} \emph{I can say that a story was broadcast, covered and brought to light, and talk about a publication and its print run.}

Complete the last lesson of chapter 286: I can say that a story was broadcast, covered and brought to light, and talk about a publication and its print run.
\end{chapteropening}

\section[\texorpdfstring{\ru{транслировать} (translírovat)}{транслировать (translírovat)}]{\ru{транслировать} (translírovat) — to broadcast}

\label{lesson:RU-C286-translirovat}



\begin{quote}
Before the new one: say the Russian for a summary, then the Russian for the result.
\end{quote}

\subsection*{You'll want to know: \ru{транслировать}}
\textbf{\ru{транслировать}} (\emph{translírovat}) — \textquotedblleft{}to broadcast\textquotedblright{}.

\begin{quote}
\ru{Матч} \ru{транслировали} \ru{по} \ru{телевизору}.
\end{quote}

\begin{quote}
\emph{Match translírovali po televízoru.}
\end{quote}

\begin{quote}
The match was broadcast on television.
\end{quote}

One form for both aspects. \emph{\ru{Трансляция}} — a broadcast.

\begin{cousinweb}
\emph{\ru{транслировать}} is usually traced to Latin \emph{trānslātus} \textquotedblleft{}carried across\textquotedblright{}, the participle of \emph{trānsferre} \textquotedblleft{}to carry over\textquotedblright{} (Lewis \& Short), with the Russian ending \emph{-\ru{ировать}}; EtymDB lists no parent.
\end{cousinweb}

\begin{grammarlens}[title={What you've built}]
One more word for telling, explaining and describing at B1.
\end{grammarlens}

\subsection*{Your turn}
\begin{itemize}
\raggedright
  \item \emph{Say it:} \emph{translírovat}
  \item \emph{Say it:} \emph{translírovat}, once more
  \item \emph{Say it:} a sentence of your own with \emph{translírovat}
  \item \emph{Recall:} say the Russian for a summary, then the Russian for the result, then say \emph{translírovat} again
\end{itemize}

\begin{tcolorbox}[breakable,colback=teal!4,colframe=teal!35!black,title={Before you move on}]
What does \ru{транслировать} mean? (\textquotedblleft{}To broadcast\textquotedblright{}.) Say it once more.
\end{tcolorbox}
\section[\texorpdfstring{\ru{освещать} (osveshchát)}{освещать (osveshchát)}]{\ru{освещать} (osveshchát) — to report on, to cover (in the news); to light up}

\label{lesson:RU-C286-osveshchat}



\begin{quote}
Before the new one: say the Russian for the result, then the Russian for to broadcast.
\end{quote}

\subsection*{You'll want to know: \ru{освещать}}
\textbf{\ru{освещать}} (\emph{osveshchát}) — \textquotedblleft{}to report on, to cover (in the news); to light up\textquotedblright{}.

\begin{quote}
\ru{Газеты} \ru{широко} \ru{освещали} \ru{эту} \ru{историю}.
\end{quote}

\begin{quote}
\emph{Gazéty shirokó osveshcháli étu istóriyu.}
\end{quote}

\begin{quote}
The newspapers covered this story widely.
\end{quote}

Imperfective; perfective \emph{\ru{осветить}} (\emph{osvetít}), \emph{\ru{я} \ru{освещу}}. First sense: \emph{\ru{лампа} \ru{освещает} \ru{комнату}} — the lamp lights the room.

\begin{cousinweb}
\emph{\ru{освещать}} is built on \emph{\ru{о}-} and \emph{\ru{свет}} \textquotedblleft{}light\textquotedblright{}, an inherited Slavic word (Proto-Slavic \emph{světŭ}), with \emph{-\ru{щ}-}, the Church Slavonic shape \emph{-\ru{свещ}-} (native Russian has \emph{\ru{ч}}, as in \emph{\ru{свеча}}): to throw light on.
\end{cousinweb}

\begin{grammarlens}[title={What you've built}]
One more word for telling, explaining and describing at B1.
\end{grammarlens}

\subsection*{Your turn}
\begin{itemize}
\raggedright
  \item \emph{Say it:} \emph{osveshchát}
  \item \emph{Say it:} \emph{osveshchát}, once more
  \item \emph{Say it:} a sentence of your own with \emph{osveshchát}
  \item \emph{Recall:} say the Russian for the result, then the Russian for to broadcast, then say \emph{osveshchát} again
\end{itemize}

\begin{tcolorbox}[breakable,colback=teal!4,colframe=teal!35!black,title={Before you move on}]
What does \ru{освещать} mean? (\textquotedblleft{}To report on, to cover (in the news); to light up\textquotedblright{}.) Say it once more.
\end{tcolorbox}
\section[\texorpdfstring{\ru{раскрывать} (raskryvát)}{раскрывать (raskryvát)}]{\ru{раскрывать} (raskryvát) — to reveal, to uncover; to solve (a case)}

\label{lesson:RU-C286-raskryvat}



\begin{quote}
Before the new one: say the Russian for to broadcast, then the Russian for to report on.
\end{quote}

\subsection*{You'll want to know: \ru{раскрывать}}
\textbf{\ru{раскрывать}} (\emph{raskryvát}) — \textquotedblleft{}to reveal, to uncover; to solve (a case)\textquotedblright{}.

\begin{quote}
\ru{Журналист} \ru{раскрыл} \ru{тайну} \ru{компании}.
\end{quote}

\begin{quote}
\emph{Zhurnalíst raskrýl táynu kompánii.}
\end{quote}

\begin{quote}
The journalist uncovered the company's secret.
\end{quote}

Imperfective; perfective \emph{\ru{раскрыть}} (\emph{raskrýt}), \emph{\ru{я} \ru{раскрою}}. \emph{\ru{Раскрыть} \ru{тайну}} — to reveal a secret.

\begin{cousinweb}
\emph{\ru{раскрывать}} is built on \emph{\ru{рас}-} \textquotedblleft{}apart\textquotedblright{} and \emph{\ru{крыть}} \textquotedblleft{}to cover\textquotedblright{}, an inherited Slavic verb: to take the cover off. EtymDB lists no parent for the compound.
\end{cousinweb}

\begin{grammarlens}[title={What you've built}]
One more word for telling, explaining and describing at B1.
\end{grammarlens}

\subsection*{Your turn}
\begin{itemize}
\raggedright
  \item \emph{Say it:} \emph{raskryvát}
  \item \emph{Say it:} \emph{raskryvát}, once more
  \item \emph{Say it:} a sentence of your own with \emph{raskryvát}
  \item \emph{Recall:} say the Russian for to broadcast, then the Russian for to report on, then say \emph{raskryvát} again
\end{itemize}

\begin{tcolorbox}[breakable,colback=teal!4,colframe=teal!35!black,title={Before you move on}]
What does \ru{раскрывать} mean? (\textquotedblleft{}To reveal, to uncover; to solve (a case)\textquotedblright{}.) Say it once more.
\end{tcolorbox}
\section[\texorpdfstring{\ru{публикация} (publikátsiya)}{публикация (publikátsiya)}]{\ru{публикация} (publikátsiya) — a publication}

\label{lesson:RU-C286-publikatsiya}



\begin{quote}
Before the new one: say the Russian for to report on, then the Russian for to reveal.
\end{quote}

\subsection*{You'll want to know: \ru{публикация}}
\textbf{\ru{публикация}} (\emph{publikátsiya}) — \textquotedblleft{}a publication\textquotedblright{}.

\begin{quote}
\ru{После} \ru{публикации} \ru{книга} \ru{стала} \ru{известной}.
\end{quote}

\begin{quote}
\emph{Pósle publikátsii kníga stála izvéstnoy.}
\end{quote}

\begin{quote}
After publication the book became well known.
\end{quote}

Feminine.

\begin{cousinweb}
\emph{\ru{публикация}} goes back to Latin \emph{pūblicātiō} \textquotedblleft{}a making public\textquotedblright{}, from \emph{pūblicus} \textquotedblleft{}of the people\textquotedblright{} (Lewis \& Short); EtymDB lists no parent for the Russian word.
\end{cousinweb}

\begin{grammarlens}[title={What you've built}]
One more word for telling, explaining and describing at B1.
\end{grammarlens}

\subsection*{Your turn}
\begin{itemize}
\raggedright
  \item \emph{Say it:} \emph{publikátsiya}
  \item \emph{Say it:} \emph{publikátsiya}, once more
  \item \emph{Say it:} a sentence of your own with \emph{publikátsiya}
  \item \emph{Recall:} say the Russian for to report on, then the Russian for to reveal, then say \emph{publikátsiya} again
\end{itemize}

\begin{tcolorbox}[breakable,colback=teal!4,colframe=teal!35!black,title={Before you move on}]
What does \ru{публикация} mean? (\textquotedblleft{}A publication\textquotedblright{}.) Say it once more.
\end{tcolorbox}
\section[\texorpdfstring{\ru{тираж} (tirázh)}{тираж (tirázh)}]{\ru{тираж} (tirázh) — a print run, a circulation; a lottery draw}

\label{lesson:RU-C286-tirazh}



\begin{quote}
Before the new one: say the Russian for to reveal, then the Russian for a publication.
\end{quote}

\subsection*{You'll want to know: \ru{тираж}}
\textbf{\ru{тираж}} (\emph{tirázh}) — \textquotedblleft{}a print run, a circulation; a lottery draw\textquotedblright{}.

\begin{quote}
\ru{У} \ru{этой} \ru{газеты} \ru{большой} \ru{тираж}.
\end{quote}

\begin{quote}
\emph{U étoy gazéty bolshóy tirázh.}
\end{quote}

\begin{quote}
This newspaper has a large circulation.
\end{quote}

Masculine; the stress moves to the ending: \emph{\ru{тиража}}, \emph{\ru{тиражи}}. \emph{\ru{Выиграть} \ru{в} \ru{тираже}} — to win in a lottery draw.

\begin{cousinweb}
\emph{\ru{тираж}} was borrowed from French \emph{tirage} \textquotedblleft{}a printing; a drawing of lots\textquotedblright{} (EtymDB). The French noun is usually traced to \emph{tirer} \textquotedblleft{}to draw, to pull\textquotedblright{}, as copies are drawn off a press and numbers out of a drum; EtymDB lists no parent for \emph{tirage} itself.
\end{cousinweb}

\begin{grammarlens}[title={What you've built}]
One more word for telling, explaining and describing at B1.
\end{grammarlens}

\subsection*{Your turn}
\begin{itemize}
\raggedright
  \item \emph{Say it:} \emph{tirázh}
  \item \emph{Say it:} \emph{tirázh}, once more
  \item \emph{Say it:} a sentence of your own with \emph{tirázh}
  \item \emph{Recall:} say the Russian for to reveal, then the Russian for a publication, then say \emph{tirázh} again
\end{itemize}

\begin{tcolorbox}[breakable,colback=teal!4,colframe=teal!35!black,title={Before you move on}]
What does \ru{тираж} mean? (\textquotedblleft{}A print run, a circulation; a lottery draw\textquotedblright{}.) Say it once more.
\end{tcolorbox}
\section[(dialogue)]{Words from chapters 279–286 — a review}

\label{lesson:RU-R286-review-279-286}



\begin{quote}
Say the Russian for a piece of news and a bulletin.
\end{quote}

\subsection*{Your turn}
Say these aloud:

\begin{itemize}
\raggedright
  \item a piece of news, a bulletin, a headline, an article, a sensation
  \item to report, to deny, to publish, to confirm, to refute
  \item a correspondent, a report, a presenter, the air, a TV channel
  \item a source, an eyewitness, a rumour, a detail, a fact
  \item to claim, to mention, to clarify, to find out, to warn
  \item an interview, a survey, an editor, to comment on, to quote
  \item an event, an incident, a news round-up, a summary, the result
  \item to broadcast, to report on, to reveal, a publication, a print run
\end{itemize}

\begin{tcolorbox}[breakable,colback=teal!4,colframe=teal!35!black,title={Before you move on}]
A piece of news? (\textbf{\ru{известие}}.) A bulletin? (\textbf{\ru{выпуск}}.) A headline? (\textbf{\ru{заголовок}}.) An article? (\textbf{\ru{статья}}.) A sensation? (\textbf{\ru{сенсация}}.) To report? (\textbf{\ru{сообщать}}.) To deny? (\textbf{\ru{отрицать}}.) To publish? (\textbf{\ru{публиковать}}.) To confirm? (\textbf{\ru{подтверждать}}.) To refute? (\textbf{\ru{опровергать}}.) A correspondent? (\textbf{\ru{корреспондент}}.) A report? (\textbf{\ru{репортаж}}.) A presenter? (\textbf{\ru{ведущий}}.) The air? (\textbf{\ru{эфир}}.) A TV channel? (\textbf{\ru{телеканал}}.) A source? (\textbf{\ru{источник}}.) An eyewitness? (\textbf{\ru{очевидец}}.) A rumour? (\textbf{\ru{слух}}.) A detail? (\textbf{\ru{подробность}}.) A fact? (\textbf{\ru{факт}}.) To claim? (\textbf{\ru{утверждать}}.) To mention? (\textbf{\ru{упоминать}}.) To clarify? (\textbf{\ru{уточнять}}.) To find out? (\textbf{\ru{выяснять}}.) To warn? (\textbf{\ru{предупреждать}}.) An interview? (\textbf{\ru{интервью}}.) A survey? (\textbf{\ru{опрос}}.) An editor? (\textbf{\ru{редактор}}.) To comment on? (\textbf{\ru{комментировать}}.) To quote? (\textbf{\ru{цитировать}}.) An event? (\textbf{\ru{событие}}.) An incident? (\textbf{\ru{происшествие}}.) A news round-up? (\textbf{\ru{хроника}}.) A summary? (\textbf{\ru{сводка}}.) The result? (\textbf{\ru{итог}}.) To broadcast? (\textbf{\ru{транслировать}}.) To report on? (\textbf{\ru{освещать}}.) To reveal? (\textbf{\ru{раскрывать}}.) A publication? (\textbf{\ru{публикация}}.) A print run? (\textbf{\ru{тираж}}.)
\end{tcolorbox}
